\subsection{Cuts and breakdown}\label{app:cuts}

\Cref{prop:control} allows one ultimatum and a cut that lasts one period. A threat of permanent exclusion withdraws instruments at many dates, and \cref{cor:windows} lets Washington issue ultimatums at several dates. The lemma below gives the bite and the cost of such a cut. Its two scope conditions push in opposite directions. If a cut drives Europe out at some date, Europe's loss exceeds the formula. If Europe responds to a cut it learns of in advance, for instance with the autonomy investment of \cref{sec:controlsetup}, the loss falls short of it.

\begin{lemma}[Spent and self-limiting cuts]\label{lem:selflimit}
Let dates $s=1,\dots,N$ have premia $Q_s$, slack $\Lambda_s=\pos{-Q_s}$ and discount factor $\delta$. A cut over a set $S$ of dates withdraws at each $s\in S$, before the date-$s$ contest, a bundle with unconditional part $u_s$, conditional part $d_s$ and direct ex post cost $g_s\ge 0$ to the United States. Suppose that at every $s\in S$ the United States still retains Europe after the cut ($Q_s+d_s<V_s$ and $Q_s+d_s\le\bar m_U$), and that the cut does not change the premium at any other date. Then:
\begin{enumerate}[label=(\roman*)]
\item the bite of the cut, Europe's loss net of the rents the United States must then pay, is $\sum_{s\in S}\delta^{s-1}\big(u_s+\min\{d_s,\Lambda_s\}\big)$;
\item the cost to the United States of carrying it out is $\sum_{s\in S}\delta^{s-1}e_s$ with $e_s=g_s+\pos{Q_s+d_s}-\pos{Q_s}$; with a separate stake at each date, the date-$s$ withdrawal is carried out if and only if $e_s\le m_s\le\bar m_U$;
\item extending a cut into a date with $Q_s\ge 0$ adds $\delta^{s-1}u_s$ to its bite and $\delta^{s-1}(g_s+d_s)$ to its cost to the United States. A purely conditional instrument ($u_s=0$) extended into such a date adds nothing to the bite.
\end{enumerate}
\end{lemma}

Part (iii) governs threats of permanent exclusion. Whether such a threat is credible depends on how much Washington can stake, but a longer cut adds no bite once Europe's premium is positive. At those dates a conditional cut hands Europe exit as its alternative, and Washington pays the difference as rent.

Under complete information an ultimatum is accepted or never issued (\cref{prop:control}). Breakdown, a refusal followed by the executed threat, requires that Washington be uncertain about what a transfer would cost Europe. The proposition below lets the sovereignty cost $\sigma$ be Europe's private information and lets Washington choose between a threat-backed ultimatum and an offer to buy control without staking the threat. We keep the other assumptions of \cref{prop:control} and, to isolate the screening problem, set $\Delta_c=\infty$, rule out the autonomy investment, and assume that Europe's period-1 premium is negative and that the threat is credible. We write $F'$ for the density of $\sigma$ and $\Pi_U$ for Washington's objective, to keep them apart from the firm penalty $f$ of \cref{eq:Q} and the objective of \cref{prop:conditionality}.

\begin{proposition}[Breakdown]\label{prop:breakdown}
Assume the conditions of \cref{prop:control} with $\Delta_c=\infty$, no investment, $Q_1<0$ and $e\le\bar m_U$, and let $\sigma$ be Europe's private information, distributed with CDF $F$ on $[0,\bar\sigma]$, density $F'>0$ on $(0,\bar\sigma)$, and $F$ log-concave. Let $R=\delta\pos{Q_2}$, $\tilde b=B_1-r^c$ and $\Pi_U(s;A,x)=F(s)(A-s)-(1-F(s))x$. For $A+x>0$ let $s^*(A,x)$ be the unique solution in $(0,\bar\sigma)$ of $A+x-s=F(s)/F'(s)$ if one exists and $\bar\sigma$ otherwise; let $s^*(A,x)=0$ if $A+x\le 0$.
\begin{enumerate}[label=(\alph*)]
\item \emph{Threat-backed ultimatum.} With payment $T$, Europe accepts if and only if $\sigma\le T+\tilde b-R$. The optimal cutoff for the United States is $\hat s=\max\{\tilde b-R,\,\min\{\bar\sigma,s^*\}\}$ with $s^*=s^*(\tilde b+v_c,e)$, its payment is $\hat s-\tilde b+R\ge 0$, and its gain is $\Pi_B=\Pi_U(\hat s;\tilde b+v_c,e)-\iota$. Conditional on this ultimatum, Europe refuses and the threat is carried out with probability $P_b=1-F(\hat s)$. For uniform $F$, $s^*=\min\{\bar\sigma,\,(\tilde b+v_c+e)/2\}$.
\item \emph{Purchase offer.} If the United States instead offers to buy control without staking the threat, Europe accepts if and only if $\sigma\le T-r^c-R$, the optimal cutoff is $s_P=s^*(v_c-r^c,0)$, and the gain is $\Pi_P=\Pi_U(s_P;v_c-r^c,0)-\iota$. If $v_c\le r^c$, then $\Pi_P\le -\iota$.
\item The United States makes an offer if and only if $\max\{\Pi_B,\Pi_P\}>0$ and backs it with the threat if and only if $\Pi_B\ge\Pi_P$. If $v_c\le r^c$ (in particular if $v_c=0$), it issues if and only if $F(\hat s)(\tilde b+v_c-\hat s)-(1-F(\hat s))e>\iota$, and every offer is threat-backed.
\item Across parameter values at which the United States issues a threat-backed ultimatum, $P_b$ is nonincreasing in $B_1$, $e$ and $v_c$, nondecreasing in $r^c$ and $R$, and independent of $\iota$. These comparative statics hold for every continuous $F$ (not necessarily log-concave) when $\hat s$ is the largest, or the smallest, maximizer.
\item If $v_c>r^c$ and the United States chooses between (a) and (b), then across parameter values at which it makes an offer the refusal probability is nonincreasing in $B_1$ and $v_c$, and the breakdown probability (refusal followed by the executed threat) is nonincreasing in $e$ and nondecreasing in $R$. The refusal probability is nondecreasing in $r^c$ on any range where the threat-backed payment is positive. The other directions can fail: the refusal probability can rise with $e$ and fall with $R$, and the breakdown probability can rise with $B_1$ and $v_c$ and fall with $r^c$.
\end{enumerate}
\end{proposition}

Given a threat-backed ultimatum, breakdown is more likely when China's punishment of a transfer is credible and large, when Europe's future premium is high, and when the threat is cheap for Washington to carry out. The last is the risk-return trade-off of \citet{fearon1995rationalist} and \citet{powell1999shadow}: a patron that loses little from a refusal makes a tougher offer and is refused more often. A larger bite makes refusal costlier for Europe, and a higher value of control makes Washington willing to pay more, so both lower the chance of refusal. When $v_c>r^c$, a patron whose threat is expensive to carry out may prefer to offer money alone, since a refused purchase offer costs it nothing, and a rise in the bite can then move it to the threat and raise the probability of breakdown (part (e)). A credible Chinese punishment is now carried out on the path with positive probability, because the types that accept are punished.

The result rests on a one-shot ultimatum, which is a commitment assumption. A patron that could not commit to walk away after a refusal, and could renew its offer quickly, would compete with its own later offers. With one-sided private information and offers made by the uninformed side, that side loses most of its power to screen as the time between offers shrinks \citep{fudenberg1985infinite, gul1986foundations}. We expect such a patron to extract less and delay to replace some of the refusals that \cref{prop:breakdown} counts as breakdown. Each refusal would also execute a cut, which brings the repeated game closer to a war of attrition; we do not solve it. The probabilities in \cref{prop:breakdown} are those of a patron that can commit to one public demand, a commitment that a staked promise to make no further offer could supply (\cref{sec:slackprinciple}).
